\subsection{Regular algebras}

DEFINITION 1. --- Let K be a field. An algebra A over K is said to be regular if L ⊗K A is an integral domain for every extension L of K.

A regular algebra is in particular an integral domain, hence commutative.

PROPOSITION 2. --- Let A and B be two algebras over a field K. If A is an integral domain and B is regular then \(A \otimes_{K} B\) is an integral domain.

Let \(L\) be the field of fractions of \(A\) . Since \(B\) is a regular \(K\) -algebra, \(L \otimes_K B\) is an integral domain, hence so is \(A \otimes_K B\) , since it is isomorphic to a subring of \(L \otimes_K B\) .

PROPOSITION 3. --- Let K be a field.

\begin{enumerate}
\def\labelenumi{\alph{enumi})}
\item
  Every subalgebra of a regular K-algebra is itself regular.
\item
  The tensor product of two regular K-algebras is again regular.
\item
  Let A be a K-algebra and \(K'\) an extension of K. For A to be regular it is necessary and sufficient that the \(K'\) -algebra \(\mathrm{A}_{(\mathbb{K}')}\) derived from A by extension of scalars should be regular.
\end{enumerate}

The proof of \(a\) (resp. \(c\) ) is identical to that of part \(a\) (resp. \(d\) ) of Prop. 3, V, p.~119, after replacing everywhere « reduced ring » by « integral domain » and « separable algebra » by « regular algebra ». Let us prove \(b\) ).

Let A and B be two regular K-algebras. Let L be an extension of K. Since A is regular, the ring \(L \otimes_{K} A\) is an integral domain. By Prop. 2, the ring \((L \otimes_{K} A) \otimes_{K} B\) is an integral domain, because B is regular. Finally, the ring \(L \otimes_{K} (A \otimes_{K} B)\) is isomorphic to \((L \otimes_{K} A) \otimes_{K} B\) , and hence is an integral domain. This shows \(A \otimes_{K} B\) to be a regular K-algebra.

